%% sections/09-parameters.tex
%% Parameters & Settings Reference
%% IEC/IEEE 82079-1:2019 — Reference Information (§5)
%% ─────────────────────────────────────────────────────

\section{Parameters and Settings Reference}
\label{sec:parameters}

%%
%% WRITING GUIDE — PARAMETERS:
%%
%%   List EVERY parameter touched in the procedure. Operators reference
%%   this section when making future adjustments without re-reading the
%%   full procedure.
%%
%%   Use the \begin{paramtable} environment.
%%   Columns: Parameter name | Description | Location | Valid Range | Unit | Default
%%
%%   For ACS parameters: Location = "ACS axis N" or "ACSPL+ global"
%%   For Python variables: Location = "force_logger.py" or similar
%%

\subsection{ACS SPiiPlus Controller Parameters}
\label{sec:params:acs}

\begin{paramtable}

  \paramrow{KP(0)}
    {Proportional gain — force PID}
    {ACS Axis 0}
    {0.1 – 20.0}
    {\si{\ampere\per\newton}}
    {2.5}

  \paramrow{KI(0)}
    {Integral gain — force PID}
    {ACS Axis 0}
    {0.0 – 5.0}
    {\si{\ampere\per\newton\per\second}}
    {0.08}

  \paramrow{KD(0)}
    {Derivative gain — force PID}
    {ACS Axis 0}
    {0.0 – 0.1}
    {\si{\ampere\second\per\newton}}
    {0.001}

  \paramrow{DCURR(0)}
    {Continuous current limit (safety ceiling)}
    {ACS Axis 0}
    {0.1 – \textit{motor spec}}
    {\si{\ampere}}
    {1.5}

  \paramrow{ECURR(0)}
    {Emergency current limit}
    {ACS Axis 0}
    {DCURR – \textit{motor spec}}
    {\si{\ampere}}
    {2.0}

  \paramrow{XVEL(0)}
    {Velocity limit during force mode}
    {ACS Axis 0}
    {0.0 – 100.0}
    {\si{\milli\metre\per\second}}
    {5.0}

  \paramrow{KP\_FFW(0)}
    {Velocity feedforward gain}
    {ACS Axis 0}
    {0.0 – 1.0}
    {---}
    {0.0}

\end{paramtable}


\subsection{ACSPL+ Global Variables}
\label{sec:params:acspl}

\begin{paramtable}

  \paramrow{force\_offset}
    {Sensor zero-offset voltage (set at zeroing step)}
    {ACSPL+ global}
    {–0.5 to +0.5}
    {\si{\volt}}
    {0.0}

  \paramrow{SENSOR\_SCALE}
    {Force sensor sensitivity (Newton per Volt)}
    {ACSPL+ global}
    {10.0 – 500.0}
    {\si{\newton\per\volt}}
    {100.0}

  \paramrow{force\_setpoint}
    {Target force applied to workpiece}
    {ACSPL+ global}
    {0.0 – 200.0}
    {\si{\newton}}
    {5.0}

\end{paramtable}


\subsection{Python Script Variables}
\label{sec:params:python}

The following variables are set at the top of \filepath{code/force\_logger.py}:

\begin{configcode}[force\_logger.py — configurable variables]
\begin{lstlisting}
# Network
ACS_IP       = "10.0.0.100"   # ACS controller IP address
ACS_PORT     = 701            # Default ACS Ethernet port (do not change)

# Logging
SAMPLE_RATE  = 100            # [Hz]  — samples per second
DURATION     = 30             # [s]   — total logging duration
OUTPUT_FILE  = "force_log.csv"

# Sensor channel (0-indexed, AIN 1 = index 0)
AIN_CHANNEL  = 0

# Calibration (must match SENSOR_SCALE in ACSPL+)
SENSOR_SCALE = 100.0          # [N/V]
FORCE_OFFSET = 0.0            # [V]   — updated at runtime from ACSPL+ global
\end{lstlisting}
\end{configcode}


\subsection{File Locations Reference}
\label{sec:params:files}

\begin{center}
\begin{tabularx}{\textwidth}{@{} l X @{}}
  \toprule
  \textbf{File} & \textbf{Location / Purpose} \\
  \midrule
  \filepath{GS500\_ForceCtrl\_baseline.cfg} &
    Request from PBA Systems support.
    ACS controller baseline configuration for GS-500 force control. \\
  \rowcolor{PBALightGray}
  \filepath{code/force\_logger.py} &
    Python script to log force data via \code{acspy} library.
    See \cref{app:code-python}. \\
  \filepath{code/plot\_force\_results.py} &
    Python script to generate acceptance plots from CSV log.
    See \cref{app:code-python}. \\
  \rowcolor{PBALightGray}
  \filepath{code/force\_control\_loop.acspl} &
    ACSPL+ force control program (loaded to ACS buffer 5).
    See \cref{app:code-acs}. \\
  \bottomrule
\end{tabularx}
\end{center}
